\section{Discussion}
\label{sec:discussion}

VED proposes that the foundational primitives of physics ---
spacetime, fields, and symmetry groups --- are not inputs but outputs.
They emerge from the dynamics of a self-referential selection process
operating on a discrete causal network.
The present work should be understood as a structural proposal, not a finalized physical theory.
Several aspects of this proposal merit explicit discussion.

\paragraph{Relation to existing frameworks.}
VED does not quantize general relativity, nor does it discretize
quantum field theory.
Its closest structural relatives are causal set theory
\cite{bombelli1987causal,sorkin2005causal},
in which spacetime is replaced by a partial order on discrete events,
and relational quantum mechanics \cite{rovelli1996relational},
which locates physical quantities in relations between systems rather
than in absolute values.
VED differs from both in that it derives spatial and temporal structure
from a single self-referential equation, without postulating an
underlying causal order.

\paragraph{Status of conservation laws.}
The continuity equation~\cref{eq:continuity} does not impose conservation.
Conservation $(\sigma = 0)$ is a special case that holds in the
classical-active phase.
This is a departure from standard formulations, where conservation is
an axiom or a consequence of symmetry (Noether's theorem).
In VED, Noether's theorem applies within the classical phase as an
effective statement; it does not hold in general.

\paragraph{The Lagrangian as an effective description.}
An early version of this framework attempted a Lagrangian formulation.
This was abandoned for two reasons:
(i) the Lagrangian principle $\delta S = 0$ presupposes a reversible
time parameter $t$, but in VED, time $\tau$ must be derived, not assumed;
(ii) the non-zero log-generation rate $\sigma$ is structurally
incompatible with a variational principle that requires conservation.
The Lagrangian formalism is, however, recoverable as an effective
description within the classical-active phase, where $\sigma \approx 0$.

\paragraph{Open problems.}
The framework presented here is structurally complete but
quantitatively open.
Key open problems include:
\begin{enumerate}
  \item Derivation of the specific value of $\lambda$ from first principles.
  \item Proof of existence and uniqueness of fixed points for generic $H_{ij}$.
  \item Quantitative derivation of mass ratios from $\Delta_k$ (Vol.~2).
  \item Exact (non-linear) Schwarzschild solution from the fixed-point
    equation (Vol.~3).
  \item Microscopic computation of Newton's constant $G$ from
    $\lambda$, $\kappa_B$, and network density (Vol.~3).
\end{enumerate}
